\begin{proof}
Put
\[
 d_\infty(\lVert\cdot\rVert_1,\lVert\cdot\rVert_2)
 :=\sup_C\left|\log
       \frac{\lVert\cdot\rVert_1}{\lVert\cdot\rVert_2}\right|.
\]
The quotient in this formula is independent of the local frame.

\emph{Step 1: passage from CLD weights to a global Thuillier metric.}
CLD-psh functions are classically psh by
\cite[\S1.2]{GR25Classical}; on a smooth analytic curve classical psh and
Thuillier-subharmonic coincide by
\cite[\S4.1.2]{GR25Classical}. In the present algebraically closed,
nontrivially valued scope one can equivalently use
\cite[Thm.~4.7 and Cor.~4.8]{Wanner19}: positivity of
\(\dprimeddprime[u]\) for a continuous function is equivalent to
Thuillier subharmonicity.

This assertion is compatible with the metric cocycle, not merely with one
chosen frame. If \(s=f s'\) for an invertible analytic function \(f\), then
\[
 -\log\lVert s\rVert
 =-\log\lVert s'\rVert-\log|f|.
\]
The function \(\log|f|\) is harmonic by
\cite[Prop.~2.3.20]{Thu05}. Thus the local subharmonic weights satisfy
exactly \cite[Def.~3.4.17]{Thu05} and define a global subharmonic metric on
\(L\). Its curvature is a positive Radon measure. The
Poincar\'e--Lelong formula \cite[Prop.~3.3.15]{Thu05} gives
\[
             \int_C c_1(L,\lVert\cdot\rVert)=\deg L,
\]
so \(\deg L\geq0\).

\emph{Step 2: Thuillier's global lisse approximation and uniformity.}
By \cite[Thm.~3.4.19]{Thu05}, the metric is the upper envelope of a
directed increasing family
\((\lVert\cdot\rVert_\alpha)_\alpha\) of global lisse subharmonic metrics
on \(L\). The occurrence of an undefined open set in the printed theorem
does not restrict this assertion: its proof constructs
\(\lVert\cdot\rVert_S\) globally on \(L\), checks the transition functions,
and states that their upper envelope is the original metric.

Set
\[
 q_\alpha:=\log
 \frac{\lVert\cdot\rVert}{\lVert\cdot\rVert_\alpha}.
\]
The \(q_\alpha\) are continuous nonnegative functions on the compact space
\(C\), decrease as a directed family, and have pointwise infimum zero.
The net form of Dini's theorem gives
\(\lVert q_\alpha\rVert_\infty\to0\): for \(\varepsilon>0\), the
increasing open sets \(\{q_\alpha<\varepsilon\}\) cover \(C\); a finite
subcover and directedness give one index whose open set is all of \(C\).
Choosing one index for each accuracy \(2^{-n}\) produces a sequence. It
remains to approximate one global lisse subharmonic metric.

\emph{Step 3: removal of type-III corners.}
Fix such a lisse subharmonic metric
\(\lVert\cdot\rVert_\ell\). Choose a formal reference metric
\(\lVert\cdot\rVert_0\) on \(L\), which exists by
\cite[\S2.2]{GM19}. Formal metrics are PL by
\cite[Prop.~2.10]{GM19}, hence their local weights are lisse by
\cite[Prop.~3.2.5]{Thu05}. We choose a representing model with reduced
special fibre, as in the first paragraph of the proof of
\cite[Prop.~5.1]{GR25PL}; thus its curvature is supported at finitely many
type-II divisorial points. The relative weight is
lisse. By \cite[Prop.~3.2.4]{Thu05} it factors through a finite semistable
skeleton. After a common semistable refinement, using
\cite[Lem.~2.2.22]{Thu05}, we may assume that this skeleton also contains
the finite type-II support of the curvature of the reference metric. Thus
we may write
\[
 -\log\frac{\lVert\cdot\rVert_\ell}
                 {\lVert\cdot\rVert_0}=\Phi\circ\tau_\Gamma
\]
on a finite semistable skeleton \(\Gamma\); no field extension is needed
because \(K\) is algebraically closed. The curvature \(\mu_\ell\) is a finite positive
atomic measure supported at points of type II or III. On \(\Gamma\) one has
\[
                 \mu_\ell=\mu_0+\Delta_\Gamma\Phi,
\]
where \(\Delta_\Gamma\) denotes the graph Laplacian with the
curvature-change sign convention in this identity.

For every type-III atom \(x\) of mass \(m_x\), choose a type-II point
\(y_x\) on the same edge, as close to \(x\) as desired, and put
\[
 \mu':=\mu_\ell-
       \sum_xm_x\delta_x+\sum_xm_x\delta_{y_x}.
\]
The value group of an algebraically closed nontrivially valued field is
divisible and nonzero, hence its points are dense on every skeletal edge;
the \(y_x\) can therefore be chosen as claimed. The measure \(\mu'\) is
positive and has the same total mass as \(\mu_\ell\).

On a compact metric graph a signed measure of total mass zero has a
piecewise affine potential, unique modulo constants. Let \(g\) be
normalized by \(g(p)=0\) and satisfy
\[
                    \Delta_\Gamma g=\mu'-\mu_\ell.
\]
The normalized graph Green kernel is continuous on
\(\Gamma\times\Gamma\). Consequently
\[
 \lVert g\rVert_\infty\longrightarrow0
\quad\text{as}\quad
 \max_x d_\Gamma(x,y_x)\longrightarrow0.
\]

Define the new metric by the global relative weight
\[
 -\log\frac{\lVert\cdot\rVert_{\ell}'}
                 {\lVert\cdot\rVert_0}
       :=(\Phi+g)\circ\tau_\Gamma.
\]
This changes no transition function, has curvature \(\mu'\), and satisfies
\(d_\infty(\lVert\cdot\rVert_{\ell}',
\lVert\cdot\rVert_\ell)=\lVert g\rVert_\infty\). It is therefore a lisse
metric arbitrarily close to \(\lVert\cdot\rVert_\ell\). All corners of its relative
skeletal weight are now at type-II points: a type-III corner would give a
type-III atom, while the formal reference has none. After subdivision at
the finitely many type-II corners, the relative weight is an
\(\mathbf R\)-PL function in the strict G-topology. This is the type-II
case of the skeletal description in
\cite[Props.~3.2.4--3.2.5]{Thu05}, written in the terminology of
\cite[\S2.5]{GR25PL}: at points of type I, III, or IV the harmonic and
\(\mathbf R\log|\mathcal O^\times|\) germs agree by
\cite[Thm.~2.3.21]{Thu05}, while at a type-II corner the G-local
generation by such logarithms is exactly
\cite[Prop.~3.2.5]{Thu05}. Since every nonharmonic corner is now type II,
the subdivision is realized by a semistable formal refinement and the
inverse images of its cells form a G-cover by strict affinoid domains. By
\cite[Lem.~2.2.22(i)]{Thu05}, no scalar extension is needed for this
type-II subdivision because \(K\) is algebraically closed.
Equivalently, the resulting function satisfies
\cite[Def.~5.1 and Cor.~5.9]{GJR26}.
This last restriction is essential: Thuillier's G-site also permits
non-strict affinoid cuts, whereas the R-PL convention in
\cite[\S2.5]{GR25PL} uses the G-topology generated by strictly analytic
domains. The warning following \cite[Cor.~5.5]{GR25PL} records precisely
the resulting obstruction for real coefficients; the explicit example is
\cite[Ex.~5.3]{GJR26}. Thus a type-III bend need not be R-PL in that strict
sense.

\emph{Step 4: the real model and its nefness.}
Because \(C\) is algebraic and proper, a global \(\mathbf R\)-PL function
is represented on one formal model by a model of the trivial bundle with
real coefficients. This is the global model passage used in the proof of
\cite[Thm.~5.6]{GR25PL}: there the analytic neighborhood is first placed
in a proper algebraic curve, and the PL-to-formal passage is
\cite[Prop.~2.10]{GM19}. Adding it to
the formal reference represents
\(\lVert\cdot\rVert_{\ell}'\) by a real line bundle
\(\mathfrak L_{\mathbf R}\) on a proper formal model
\(\mathfrak C\). As in the proof of \cite[Prop.~5.1]{GR25PL}, over the
algebraically closed field \(K\) we may pass to a model with reduced special
fibre.

By \cite[Prop.~5.1]{GR25PL}, subharmonicity of an
\(\mathbf R\)-PL weight on a smooth curve is equivalent to nefness of its
residue line bundle. Equivalently, the coefficient of the curvature
measure at the divisorial point associated to an irreducible component
\(Y\subset\mathfrak C_s\) is, up to the fixed positive normalization,
\(\deg_Y\mathfrak L_{\mathbf R}\). Since \(\mu'\geq0\),
\[
                    \deg_Y\mathfrak L_{\mathbf R}\geq0
 \qquad\text{for every irreducible component }Y.
\]
Thus \(\mathfrak L_{\mathbf R}\) is nef; compare the formal-model
criterion \cite[Lem.~7.4]{GJR26}. Notice that this step uses the
metric on \(L\), rather than trying to change local weights independently;
the transition law has already been retained throughout.

\emph{Step 5: rational approximation inside the nef face.}
On one common proper model write
\[
 \mathfrak L_{\mathbf R}
 =\mathfrak L_0+\sum_{i=1}^r a_i\mathfrak M_i,
 \qquad a_i\in\mathbf R,
\]
where \(\mathfrak L_0\) is an integral model of \(L\) and each
\(\mathfrak M_i\) is an integral vertical model of the trivial bundle.
If \(Y_1,\ldots,Y_N\) are the irreducible components of the special fibre,
put
\[
 P:=\left\{x\in\mathbf R^r:
 \deg_{Y_j}\mathfrak L_0+
 \sum_i x_i\deg_{Y_j}\mathfrak M_i\geq0
 \text{ for }1\leq j\leq N\right\}.
\]
This is a rational polyhedron and \(a=(a_i)\in P\). Rational points of a
rational polyhedron are dense at every one of its real points, including
points on the boundary: apply the assertion in the rational minimal face
containing the point and approximate in its rational affine hull. Choose
\(a^{(m)}\in P\cap\mathbf Q^r\) with \(a^{(m)}\to a\). Then
\[
 \mathfrak L_m:=\mathfrak L_0+
                 \sum_i a_i^{(m)}\mathfrak M_i
\]
is a nef \(\mathbf Q\)-model of \(L\). Clearing denominators produces an
integral nef model of a positive tensor power of \(L\). If \(h_i\) is the
model function of \(\mathfrak M_i\), then
\[
 d_\infty(\lVert\cdot\rVert_{\mathfrak L_m},
          \lVert\cdot\rVert_{\mathfrak L_{\mathbf R}})
 \leq\sum_i|a_i^{(m)}-a_i|\,\lVert h_i\rVert_\infty\longrightarrow0.
\]
This proves the required approximation of one lisse metric. Combining
Steps 2--5 and taking a diagonal sequence proves
\cref{itm:p081:model}. For a proper algebraic curve, formal, algebraic, and
PL metrics, including their semipositive versions, agree by
\cite[Thm.~1.1]{GM19}; this proves the algebraic assertion.

\emph{Step 6: degree zero and the absence of a hidden reference
assumption.}
If \(\deg L=0\), the positive curvature measure of the given metric has
total mass zero, hence vanishes. The zero-curvature characterization
\cite[\S3.12]{GR25Classical} shows that the metric is harmonic, therefore lisse,
and Steps 3--5 apply directly. Every nef rational approximant has
nonnegative curvature coefficients, while its Chambert--Loir curvature
measure has total mass \(\deg L=0\); hence all of those coefficients
vanish. In particular the proof does not assume a
semipositive degree-zero reference metric.

There is also an unconditional existence statement. By the last sentence
of \cite[Rem.~3.4.20]{Thu05}, every degree-zero line bundle on a smooth
proper analytic curve admits a harmonic metric. It is lisse; Steps 3--5
then yield a semipositive \(\mathbf Q\)-formal metric. This supplies the
degree-zero reference metric whenever one is wanted later.

\emph{Step 7: the algebraic Fubini--Study refinement in positive degree.}
Suppose that \(\deg L>0\). Then \(L\) is
ample. Given one nef rational model metric,
choose, on a model dominating its model, a relatively ample rational model
of \(L\); its existence after a tensor power is
\cite[Prop.~4.11]{GM19}. Write \(N\) for the nef rational model and \(A\)
for the ample rational model on the common dominating model. For rational
\(0<\varepsilon\ll1\), the model
\[
                         (1-\varepsilon)N+\varepsilon A
\]
has generic fibre \(L\) and represents the corresponding convex
combination of their additive weights. On the
projective special-fibre curve, nef plus a positive multiple of ample is
ample. Letting \(\varepsilon\to0\) gives relatively ample model metrics
arbitrarily close to the nef one. After a further tensor power the model
line bundle is relatively generated, and its metric is a Fubini--Study
model metric. A diagonal choice proves the final assertion.
\end{proof}

\begin{remark}
\label{rem:p081-curve-comparisons}
The theorem proves the curve case of the local-to-global question for the
nef-model/Zhang cone. In positive degree it also proves the
Fubini--Study-closure formulation. It does not assert global approximation
by ordinary CLD-smooth psh metrics in degree zero: a numerically trivial
residual line bundle need not be semiample over a general residue field.

The corresponding statement for a proper strict analytic curve which has
not first been algebraized is not claimed here. The proof of
\cite[Thm.~5.6]{GR25PL} globalizes an \(\mathbf R\)-PL function only after
placing the relevant neighborhood in a proper algebraic curve, whereas
\cite[Prop.~2.10]{GM19} treats integral PL metrics. A global real-model
gluing or algebraization argument is therefore still required for that
version of the local P08.2 question.

On continuous functions, \cite[Thm.~4.7 and Cor.~4.8]{Wanner19} and the
ordinary regularization theorem
\(\cref{thm:p063-current-psh-ordinary-regularization}\) identify positive
CLD current, Thuillier subharmonicity, and local decreasing ordinary
CLD-smooth psh approximation in the stated algebraically closed scope.
The theorem adds the global Zhang arrow while preserving the metric
cocycle. It does not use local ordinary approximants for gluing.

Finally, in the original discretely valued residue-characteristic-zero BFJ
framework, continuous BFJ psh metrics and continuous Zhang-semipositive
metrics agree by \(\cref{thm:ch21-continuous-zhang}\). This is a dictionary
statement in that separate field scope, not a base-change step in the proof
above; no field is tacitly assumed to be both algebraically closed and
discretely valued.
\end{remark}
